% META: {"id": "TCOMPL-CALCULS-AIRES-QCM", "chapitre": "TCOMPL-CALCULS-AIRES", "type_objet": "qcm", "capacites": ["TCOMPL-AIR-C1", "TCOMPL-AIR-C2", "TCOMPL-AIR-C3", "TCOMPL-AIR-C4", "TCOMPL-AIR-C5", "TCOMPL-AIR-C6"], "genere_depuis": "chapitres/TCOMPL-CALCULS-AIRES/qcm/TCOMPL-CALCULS-AIRES-QCM.json", "status": "generated"}
% Fichier genere par scripts/build_qcm_tex.py — ne pas editer a la main.

\section*{\textcolor{chapcolor}{\MakeUppercase{QCM d'auto-évaluation - Calculs d'aires}}}

\begin{center}
\textit{Pour chaque question, une seule réponse est exacte.}
\end{center}

\begin{enumerate}

\item[] \textbf{Capacité C5}

\item \textbf{[Q1]} Deux primitives d'une même fonction continue sur un intervalle différent :
  \begin{enumerate}[label=\Alph*.]
    \item d'une constante additive.
    \item d'un facteur multiplicatif.
    \item toujours de 1.
    \item elles sont necessairement égales.
  \end{enumerate}

\item[] \textbf{Capacité C4}

\item \textbf{[Q2]} Une primitive de u'/u (avec u>0) est :
  \begin{enumerate}[label=\Alph*.]
    \item ln(u)
    \item u\textasciicircum{}2
    \item e\textasciicircum{}u
    \item 1/u
  \end{enumerate}

\item[] \textbf{Capacité C1}

\item \textbf{[Q3]} Pour f continue et positive sur [a,b], l'intégrale de f de a a b représente :
  \begin{enumerate}[label=\Alph*.]
    \item la pente de la tangente a la courbe.
    \item l'aire sous la courbe entre a et b.
    \item la valeur de f en b moins celle en a.
    \item toujours un nombre négatif.
  \end{enumerate}

\item[] \textbf{Capacité C3}

\item \textbf{[Q4]} Si F est une primitive de f, alors l'intégrale de a a b de f(x)dx vaut :
  \begin{enumerate}[label=\Alph*.]
    \item F(a)-F(b)
    \item F(b)-F(a)
    \item F(a)+F(b)
    \item F(b)*F(a)
  \end{enumerate}

\item[] \textbf{Capacité C2}

\item \textbf{[Q5]} Pour une fonction croissante, la méthode des rectangles a gauche donne :
  \begin{enumerate}[label=\Alph*.]
    \item un majorant de l'intégrale.
    \item la valeur exacte de l'intégrale.
    \item un minorant de l'intégrale.
    \item une valeur négative.
  \end{enumerate}

\item[] \textbf{Capacité C6}

\item \textbf{[Q6]} Pour f continue, positive et croissante, la fonction F(x)=intégrale de a a x de f(t)dt vérifie :
  \begin{enumerate}[label=\Alph*.]
    \item F n'est pas dérivée.
    \item F=f'
    \item F est constante.
    \item F'=f
  \end{enumerate}

\end{enumerate}

% NEXUS-QCM-TEACHER-ONLY-BEGIN
\ifnxVersionProfesseur
\clearpage
\section*{Clé de correction — réservée au professeur}

\begin{center}
\begin{tabular}{lll}
\hline
Question & Capacité & Réponse exacte \\
\hline
Q1 & C5 & \textbf{A} \\
Q2 & C4 & \textbf{A} \\
Q3 & C1 & \textbf{B} \\
Q4 & C3 & \textbf{B} \\
Q5 & C2 & \textbf{C} \\
Q6 & C6 & \textbf{D} \\
\hline
\end{tabular}
\end{center}

\subsection*{Diagnostic des réponses erronées}

\begin{itemize}
  \item \textbf{Q1}
  \begin{itemize}
    \item \textbf{B} — Un facteur multiplicatif changerait la dérivée, sauf cas particulier, et ne conserve donc pas la même fonction dérivée. Comme (F-G)'=f-f=0 sur l'intervalle, deux primitives diffèrent d'une constante additive. \emph{Renvoi : C5.}
    \item \textbf{C} — La constante de différence peut être n'importe quel réel et n'est pas toujours égale à 1. Comme (F-G)'=f-f=0 sur l'intervalle, deux primitives diffèrent d'une constante additive. \emph{Renvoi : C5.}
    \item \textbf{D} — Deux primitives peuvent avoir des constantes d'intégration différentes et ne sont donc pas nécessairement égales. Comme (F-G)'=f-f=0 sur l'intervalle, deux primitives diffèrent d'une constante additive. \emph{Renvoi : C5.}
  \end{itemize}
  \item \textbf{Q2}
  \begin{itemize}
    \item \textbf{B} — La dérivée de $u^2$ vaut $2uu'$ et ne produit pas le quotient $u'/u$. La formule (ln u)'=u'/u pour u>0 montre qu'une primitive de u'/u est ln(u). \emph{Renvoi : C4.}
    \item \textbf{C} — La dérivée de $e^u$ vaut $u'e^u$, pas $u'/u$. La formule (ln u)'=u'/u pour u>0 montre qu'une primitive de u'/u est ln(u). \emph{Renvoi : C4.}
    \item \textbf{D} — La dérivée de $1/u$ vaut $-u'/u^2$, avec un signe et une puissance de dénominateur différents. La formule (ln u)'=u'/u pour u>0 montre qu'une primitive de u'/u est ln(u). \emph{Renvoi : C4.}
  \end{itemize}
  \item \textbf{Q3}
  \begin{itemize}
    \item \textbf{A} — La pente d'une tangente est une dérivée ponctuelle et non une accumulation sur tout l'intervalle. Pour f continue positive, l'intégrale entre a et b représente l'aire sous la courbe sur cet intervalle. \emph{Renvoi : C1.}
    \item \textbf{C} — La différence $f(b)-f(a)$ compare deux valeurs de la fonction et non l'aire accumulée sous sa courbe. Pour f continue positive, l'intégrale entre a et b représente l'aire sous la courbe sur cet intervalle. \emph{Renvoi : C1.}
    \item \textbf{D} — La positivité de f rend au contraire son intégrale positive ou nulle. Pour f continue positive, l'intégrale entre a et b représente l'aire sous la courbe sur cet intervalle. \emph{Renvoi : C1.}
  \end{itemize}
  \item \textbf{Q4}
  \begin{itemize}
    \item \textbf{A} — Cette différence inverse les bornes et donne l'opposé de l'intégrale demandée. Le théorème fondamental donne directement l'intégrale de f entre a et b égale à F(b)-F(a). \emph{Renvoi : C3.}
    \item \textbf{C} — Le théorème fondamental utilise une différence des valeurs de la primitive et non leur somme. Le théorème fondamental donne directement l'intégrale de f entre a et b égale à F(b)-F(a). \emph{Renvoi : C3.}
    \item \textbf{D} — Aucun produit des valeurs de la primitive n'intervient dans le calcul d'une intégrale définie. Le théorème fondamental donne directement l'intégrale de f entre a et b égale à F(b)-F(a). \emph{Renvoi : C3.}
  \end{itemize}
  \item \textbf{Q5}
  \begin{itemize}
    \item \textbf{A} — Pour une fonction croissante, la hauteur gauche est plus petite que les valeurs prises ensuite sur chaque sous-intervalle. Pour une fonction croissante, chaque hauteur prise à gauche est inférieure aux valeurs suivantes : la somme est un minorant. \emph{Renvoi : C2.}
    \item \textbf{B} — Les rectangles ne donnent la valeur exacte que dans des cas particuliers, par exemple pour une fonction constante. Pour une fonction croissante, chaque hauteur prise à gauche est inférieure aux valeurs suivantes : la somme est un minorant. \emph{Renvoi : C2.}
    \item \textbf{D} — Une fonction croissante peut être positive ou négative ; la méthode n'impose pas le signe de la somme obtenue. Pour une fonction croissante, chaque hauteur prise à gauche est inférieure aux valeurs suivantes : la somme est un minorant. \emph{Renvoi : C2.}
  \end{itemize}
  \item \textbf{Q6}
  \begin{itemize}
    \item \textbf{A} — La continuité de f est précisément l'hypothèse qui garantit la dérivabilité de la fonction intégrale. Le théorème de la fonction intégrale donne F'(x)=f(x) lorsque f est continue sur l'intervalle. \emph{Renvoi : C6.}
    \item \textbf{B} — L'égalité F=f' échange la fonction intégrale et la dérivée de l'intégrande sans justification. Le théorème de la fonction intégrale donne F'(x)=f(x) lorsque f est continue sur l'intervalle. \emph{Renvoi : C6.}
    \item \textbf{C} — Comme F'=f et que f est positive, F est croissante plutôt que constante sauf si f est nulle. Le théorème de la fonction intégrale donne F'(x)=f(x) lorsque f est continue sur l'intervalle. \emph{Renvoi : C6.}
  \end{itemize}
\end{itemize}
\fi
% NEXUS-QCM-TEACHER-ONLY-END
